% Authors: Dr. Denys Dutykh (Mathematics Department, Khalifa University of Science and Technology,
% Abu Dhabi, UAE) and Prof. Laurent Vuillon (Univ. Savoie Mont Blanc, CNRS, LAMA, Chambery, France)
\section{Complete Gaussian fibers and the finite windows}
\label{sec:gaussian}

The finite part of the proof must compare an occurrence with every
primitive representation of its label. Searching a coordinate rectangle
would obscure completeness. Gaussian factorization gives an exact finite
parameter space whose size depends on the distinct prime divisors, not
on the size of the coordinates, and once the prime factorizations are
known the enumeration is a matter of integer arithmetic alone.

\subsection{The whole primitive representation fiber}

For $M\geqslant5$, define
$$
 \mathcal F(M)=\{(x,y)\in\Z^2:0<x<y,\ \gcd(x,y)=1,\ x^2+y^2=M\}.
$$
The source construction in Section~\ref{sec:occurrences} proves that this
set is nonempty for every occurrence considered here.

\begin{proposition}[Primitive Gaussian allocation]
\label{prop:gaussian-fibre}
If $\mathcal F(M)\ne\varnothing$, then
\begin{equation}\label{eq:gaussian-factorization}
 M=2^\epsilon\prod_{j=1}^t p_j^{e_j},\qquad
 \epsilon\in\{0,1\},\quad t\geqslant1,\quad
 p_j\equiv1\pmod4,
\end{equation}
where the $p_j$ are distinct primes and $e_j\geqslant1$.
Choose $\pi_j=a_j+ib_j$ with $p_j=a_j^2+b_j^2$ and $0<a_j<b_j$.
The elements of $\mathcal F(M)$ are exactly the sorted absolute
coordinate pairs of
\begin{equation}\label{eq:gaussian-sign-cube}
 (1+i)^\epsilon\prod_{j=1}^t
 \begin{cases}
 \pi_j^{e_j},&\sigma_j=1,\\
 \overline{\pi_j}^{\,e_j},&\sigma_j=-1,
 \end{cases}
 \qquad (\sigma_1,\ldots,\sigma_t)\in\{1,-1\}^t.
\end{equation}
Each normalized pair occurs exactly twice. Equivalently, fixing
$\sigma_1=1$ lists the fiber without repetition, and
$|\mathcal F(M)|=2^{t-1}$.
\end{proposition}

\begin{proof}
The Gaussian integers form a Euclidean domain for the squared norm,
hence have unique factorization. If $z=x+iy$ is primitive and a rational
prime $p\equiv3\pmod4$ divides $z\overline z$, its Gaussian primality
forces it to divide $z$ or $\overline z$. Either condition gives
$p\mid x,y$, a contradiction. Likewise $4\nmid M$: a sum of two
squares divisible by four has both coordinates even. This proves the
shape \eqref{eq:gaussian-factorization}; $t=0$ would leave $M=1$ or $2$.

For a split prime $p_j=\pi_j\overline{\pi_j}$, an allocation using both
$\pi_j$ and $\overline{\pi_j}$ in $z$ would make $p_j$ divide both
coordinates. Primitivity therefore puts the whole exponent $e_j$ on
one of the two conjugate factors. Conversely, each allocation in
\eqref{eq:gaussian-sign-cube} has primitive coordinates: a common odd
prime would require both conjugate Gaussian factors, and a common factor
two would require a higher ramified exponent. Unique factorization proves
that these allocations, up to multiplication by a unit, exhaust all
representations. Since $M\geqslant5$, a primitive Gaussian integer of norm $M$
has nonzero coordinates of distinct absolute values, for otherwise $M=1$ or
$M=2$; its sorted absolute pair therefore lies in $\mathcal F(M)$.

It remains to count. Two Gaussian integers $x+iy$ and $x'+iy'$ have the
same normalized pair, that is $\{\abs x,\abs y\}=\{\abs{x'},\abs{y'}\}$,
exactly when $x'+iy'$ is $u(x+iy)$ or $u(x-iy)$ for one of the four units
$u$. Write $z_\sigma$ for the product \eqref{eq:gaussian-sign-cube}. For an
odd prime $p_j=a_j^2+b_j^2$ the factors $\pi_j$ and $\overline{\pi_j}$ are
not associate: $\overline{\pi_j}=u\pi_j$ would force $b_j=0$, $a_j=0$ or
$a_j=\pm b_j$, none of which is possible since $0<a_j<b_j$. Suppose now that
$z_{\sigma'}$ and $z_\sigma$ give the same normalized pair. If
$z_{\sigma'}=uz_\sigma$, unique factorization compares the exponents of the
non-associate primes $\pi_j$ and $\overline{\pi_j}$ and gives $\sigma'=\sigma$.
If $z_{\sigma'}=u\overline{z_\sigma}$, the same comparison gives
$\sigma'=-\sigma$, because conjugation exchanges $\pi_j$ and
$\overline{\pi_j}$ and replaces $1+i$ by the associate $1-i=-i(1+i)$.
Conversely $z_{-\sigma}$ is an associate of $\overline{z_\sigma}$, so $\sigma$
and $-\sigma$ always give the same pair. As $t\geqslant1$, the vectors
$\sigma$ and $-\sigma$ are distinct, so each normalized pair comes from
exactly two sign vectors, which proves the count and the last assertion.
\end{proof}

The exponents in \eqref{eq:gaussian-factorization} do not enlarge the
primitive sign cube. Mixed allocations of a prime power generate
imprimitive representations, so omitting them is a proved reduction,
not a filter based on observed costs.

\begin{remark}[Labels with a single representation]\label{rem:single-representation}
By \cref{prop:gaussian-fibre}, a label $m>2$ has exactly one primitive
positive representation precisely when $t=1$, that is when $m=p^e$ or
$m=2p^e$ for a prime $p\equiv1\pmod4$. Such a label belongs to exactly one
reduced fraction. \emph{Step 1.} Let $m=m_{P/Q}=m_{P'/Q'}$; by
\cref{prop:occurrence-cost} the two occurrence pairs are elements of
$\mathcal F(m)$, of costs $P+Q-2$ and $P'+Q'-2$. \emph{Step 2.} Since
$\abs{\mathcal F(m)}=2^{t-1}=1$, the two pairs coincide, and so do their
costs: $P+Q=P'+Q'$. \emph{Step 3.} \Cref{lem:fixed-cost-injectivity} gives
$P/Q=P'/Q'$. This recovers the theorem of Baragar, Button and Schmutz that a
Markov number which is a prime power or twice a prime power is unique
\cite{Baragar1996,Button1998,Schmutz1996}, in the form stated in
\cite[Theorem~C]{Zhang2007}; see also \cite{LangTan2007}. Here it follows
from the fixed-sum ordering \eqref{eq:fixed-sum-order} and the Gaussian
count, and \cref{thm:thirteen-rows} extends the same fiber argument, on
thirteen rows, to fibers of every size.
\end{remark}

\begin{example}[Two representations, two costs]
The first-row label $m_{1/7}=610=2\cdot5\cdot61$ has exactly two
primitive positive representations:
$$
 610=13^2+21^2=9^2+23^2.
$$
In the normalization of the proposition the two split primes are
$1+2i$ and $5+6i$, and the two allocations give
$(1+i)(1+2i)(5+6i)=-23+9i$ and $(1+i)(1-2i)(5+6i)=21+13i$. Proposition~\ref{prop:gaussian-fibre}
proves that the list is complete. Their continued fractions are
$$
 \frac{21}{13}=[1;1,1,1,1,2],\qquad
 \frac{23}{9}=[2;1,1,4],
$$
so the Euclidean costs are $6$ and $7$, respectively. The occurrence
cost $1+7-2=6$ is therefore the minimum. This illustrates why the
optimization question persists even when the label and its Markov
occurrence are fixed: the full arithmetic fiber contains another
primitive pair whose directive is unrestricted.
\end{example}

\Cref{prop:gaussian-fibre} lists the whole fiber but does not say which of
its elements comes from the occurrence. The two labels next to the
occurrence in the Markov tree do: the occurrence pair is the element of the
sign cube whose prime factors all divide one explicit Gaussian integer.

\begin{proposition}[The occurrence pair in the sign cube]\label{prop:occurrence-in-cube}
Let $0<P<Q$ be coprime, let $L=P_L/Q_L<R=P_R/Q_R$ be its two Farey
parents, the reduced fractions with $P_L+P_R=P$, $Q_L+Q_R=Q$ and
$P_RQ_L-P_LQ_R=1$, and let $m_L$, $m_R$ be their labels, with $m_{0/1}=1$ and
$m_{1/1}=2$. Write $(s,r)^{\mathsf T}=He$ for the
occurrence pair of $P/Q$ before sorting, and $m=m_{P/Q}$. Then the ideal of
$\Z[i]$ generated by $m$ and $m_L+i\,m_R$ is generated by $s+ir$ or by its
conjugate. In particular, in the sign cube \eqref{eq:gaussian-sign-cube} of
$m$, the occurrence pair is the element whose prime factors all divide
$m_L+i\,m_R$, up to the conjugation of the whole product.
\end{proposition}

\begin{proof}
\emph{The root read from the Cohn matrix.} The lower Christoffel word of
$P/Q$ is the product of those of its two Farey parents, the words of $0/1$
and $1/1$ being the letters $a$ and $b$ \cite{Reutenauer2009}. Writing
$\gamma_L$ and $\gamma_R$ for the words of the parents, we have
$\mu(\gamma_{P/Q})=\mu(\gamma_L)\,\mu(\gamma_R)$. Each factor has
determinant one and upper right entry equal to its label: for the letters,
$\mu(a)$ and $\mu(b)$ show it, and for a longer Christoffel word the
computation in the proof of \cref{prop:occurrence-cost} gives an upper right
entry $4\alpha+4\beta+\delta$, which is the label by
\eqref{eq:square-count-identity}. Writing
$\mu(\gamma_{P/Q})=\left(\begin{smallmatrix}a&m\\c&d\end{smallmatrix}\right)$ and
$\mu(\gamma_X)=\left(\begin{smallmatrix}a_X&m_X\\c_X&d_X\end{smallmatrix}\right)$
for $X=L,R$, the upper row of the product gives $m=a_Lm_R+m_Ld_R$ and
$a=a_La_R+m_Lc_R$, so that
$$
 a\,m_R=a_R(m-m_Ld_R)+m_Lc_Rm_R\equiv-m_L(a_Rd_R-m_Rc_R)=-m_L\pmod m .
$$

\emph{The root read from the occurrence pair.} By the proof of
\cref{prop:occurrence-cost}, $\mu(\gamma_{P/Q})=\mu(a)H^{\mathsf T}H\mu(b)$. With
$e_1=(1,0)^{\mathsf T}$, the symmetric matrix $\mu(a)$ satisfies $\mu(a)e_1=e$,
and $\mu(b)e_1=(5,2)^{\mathsf T}=2e+e_1$. Put $v=He=(s,r)^{\mathsf T}$ and
$h=He_1=(p,q)^{\mathsf T}$. Then
$$
 a=e_1^{\mathsf T}\mu(\gamma_{P/Q})e_1=(He)^{\mathsf T}H(2e+e_1)=2m+sp+rq .
$$
The matrix with columns $v$ and $h$ is $HP_2$, of determinant
$\eta=sq-rp=\pm1$, so that $rp=sq-\eta$ and
$$
 a\,r\equiv(sp+rq)\,r=s(sq-\eta)+r^2q=qm-\eta s\equiv-\eta s\pmod m .
$$

\emph{The ideals.} The labels $m_L$, $m_R$, $m$ form a Markov triple, whose
entries are pairwise coprime \cite{Aigner2013}; and $\gcd(s,r)=1$ with
$s^2+r^2=m$. Hence $m_R$ and $r$ are invertible modulo $m$. By the first step,
$m_L+i\,m_R\equiv m_R\,(i-a)\pmod m$, and by the second, $s\equiv-\eta ar$, so that
$s+ir\equiv r\,(i-\eta a)$ and $s-ir\equiv-r\,(i+\eta a)\pmod m$. If $\eta=1$, the
first congruence and the third give equal ideals
$(m,s+ir)=(m,i-a)=(m,m_L+i\,m_R)$; if $\eta=-1$, the fourth does,
$(m,s-ir)=(m,m_L+i\,m_R)$. Finally $(m,s\pm ir)=(s\pm ir)$, because $s\pm ir$
divides $m=(s+ir)(s-ir)$.

\emph{The sign cube.} The generator $s\pm ir$ of the ideal is the greatest
common divisor of $m$ and $m_L+i\,m_R$. A Gaussian prime factor of the label
divides it exactly when it divides $m_L+i\,m_R$, which identifies the
allocation of \eqref{eq:gaussian-sign-cube}, up to the conjugation of the
whole product.
\end{proof}

The congruence $a\,m_R\equiv-m_L\pmod m$ of the first step is classical:
up to sign, the residue of $a$ modulo $m$ is the characteristic number of
the Markov triple $(m,m_L,m_R)$, defined by Frobenius as the solution of
$m_Rx\equiv\pm m_L\pmod m$, a square root of $-1$ modulo $m$
\cite[\S2, p.~462]{Frobenius1913}, see also \cite[Section~2.2]{Aigner2013}.
The proposition translates it into the Gaussian integers and identifies the
corresponding element of the sign cube with the occurrence pair.
The cofactor $(m_L+i\,m_R)/(s\pm ir)$ has norm
$(m_L^2+m_R^2)/m=3m_Lm_R-m$, by the Markov equation, in the form
$mm'=m_L^2+m_R^2$ with $m'=3m_Lm_R-m$ that Frobenius uses
\cite[\S1, p.~460]{Frobenius1913}. For the label $610$ of
the example above, the parents $0/1$ and $1/6$ have labels $1$ and $233$, and
$$
 1+233\,i=(21+13\,i)(5+8\,i):
$$
the first factor is the occurrence pair $(1+i)(1-2i)(5+6i)$, and the second
has norm $89=3\cdot1\cdot233-610$, the label of $1/5$. The second factor is no
accident: $1/5$ is the third vertex of the Farey triangle $(0/1,1/6,1/5)$ of
$1/7$, defined below, and the next subsection shows that the cofactor is
always the occurrence pair of that vertex, up to a unit and a conjugation;
here $5+8i=i\,(8-5i)$, and $(8,5)$ is the occurrence pair of $1/5$.

\subsection{The occurrence pairs along the Farey tree}\label{sec:vieta-jump}

\Cref{prop:occurrence-in-cube} finds the occurrence pair as a greatest
common divisor, a computation that still involves the label itself. The
Markov tree offers a recursion instead. Along the tree the labels obey two
relations of Vieta, a sum and a product, and the product lifts exactly to
the Gaussian integers. The Gaussian integer $m_L+i\,m_R$ of
\cref{prop:occurrence-in-cube} is the product of two Gaussian integers
attached to $P/Q$ and to the third vertex of its Farey triangle, which are
their occurrence pairs up to units and conjugation; once an orientation is
fixed at the root of the tree, the product holds exactly, without any unit.
Every occurrence pair is thus obtained from pairs found earlier by one exact
division in $\Z[i]$, with no factorization, no greatest common divisor and
no continued fraction. This is the Gaussian form of the replacement
$(x,y,z)\mapsto(x,y,3xy-z)$ that generates the Markov tree.

Two fractions $L=P_L/Q_L<R=P_R/Q_R$ with nonnegative terms are \emph{Farey
neighbors} when $P_RQ_L-P_LQ_R=1$. Their mediant is
$M=(P_L+P_R)/(Q_L+Q_R)$, and we call their \emph{difference} the fraction $D$
whose terms are $\abs{P_R-P_L}$ and $\abs{Q_R-Q_L}$. The Farey tree of $[0,1]$
starts from the pair $(0/1,1/1)$, whose difference is $1/0$, and replaces
each pair $(L,R)$ by the two pairs $(L,M)$ and $(M,R)$. Every reduced
fraction of $(0,1)$ is the mediant of exactly one pair of the tree, the pair
of its two Farey parents \cite[Section~3.2, Theorem~3.9 and its proof]{Aigner2013}.
The differences of the two new pairs are $R$ and $L$, because $M-L=R$ and
$R-M=-L$ term by term, so that the difference of a pair of the tree is
either one of the three fractions $0/1$, $1/1$ and $1/0$, or a mediant
formed earlier. We write $\gamma_X$ and $m_X$ for the lower Christoffel word
and the label of a fraction $X$ of $[0,1]$, with $\gamma_{0/1}=a$ and
$\gamma_{1/1}=b$, we set $m_{1/0}=1$, and we call $(L,R,D)$ the \emph{Farey
triangle} of the mediant $M$. Finally, $N(g+ih)=g^2+h^2$ denotes the norm of
$\Z[i]$.

The first lemma records the two relations of Vieta on every triangle. In
this indexing they are the content of Theorem~1 of Reutenauer and of its
proof \cite[Theorem~1 and its proof]{Reutenauer2009}, where the Fricke
identity gives the Markov triple of a standard factorization and two trace
identities give the two cases below; we include the argument in the
notation of the Farey tree.

\begin{lemma}[The labels on a Farey triangle]\label{lem:farey-triangle}
For every pair $(L,R)$ of the tree, with mediant $M$ and difference $D$, the
triple $(m_L,m_R,m_D)$ is a Markov triple with pairwise coprime entries, and
\begin{equation}\label{eq:vieta-relations}
 m_M=3m_Lm_R-m_D,\qquad m_Mm_D=m_L^2+m_R^2 .
\end{equation}
\end{lemma}

\begin{proof}
\emph{The sum relation.} Put $A_X=\mu(\gamma_X)$ for every fraction $X\neq1/0$
of the tree. Since $\gamma_M=\gamma_L\gamma_R$ \cite{Reutenauer2009}, we have
$A_M=A_LA_R$, and each $A_X$ has determinant one and trace $3m_X$ by
\eqref{eq:cohn-source}. For two matrices $A$ and $B$ of determinant one, the
Cayley--Hamilton theorem gives $B+B^{-1}=(\tr B)\,I$; multiplying by $A$ and
taking traces,
$$
 \tr(AB)+\tr(AB^{-1})=\tr A\,\tr B .
$$
With $A=A_L$ and $B=A_R$ this reads $3m_M+\tr(A_LA_R^{-1})=9m_Lm_R$, and it
remains to show that $\tr(A_LA_R^{-1})=3m_D$. For the starting pair,
$$
 A_{0/1}A_{1/1}^{-1}
 =\begin{pmatrix}2&1\\1&1\end{pmatrix}\begin{pmatrix}1&-2\\-2&5\end{pmatrix}
 =\begin{pmatrix}0&1\\-1&3\end{pmatrix},
$$
whose trace is $3=3m_{1/0}$. Every other pair $(L,R)$ was created from a pair
of the tree in one of two ways. Either it is $(L,M')$, where $M'=R$ is the
mediant of a pair $(L,R')$; then $D=R'$ and $A_R=A_LA_D$, so that
$A_LA_R^{-1}=A_LA_D^{-1}A_L^{-1}$, whose trace is $\tr A_D^{-1}=\tr A_D=3m_D$, a
matrix of determinant one and its inverse having the same trace. Or it is
$(M',R)$, where $M'=L$ is the mediant of a pair $(L',R)$; then $D=L'$ and
$A_L=A_DA_R$, so that $A_LA_R^{-1}=A_D$, whose trace is $3m_D$.

\emph{The Markov triples.} The starting triple $(1,2,1)$ satisfies
$x^2+y^2+z^2=3xyz$. If $(m_L,m_R,m_D)$ satisfies it, then $m_D$ is a root of
the polynomial $X^2-3m_Lm_RX+m_L^2+m_R^2$, and by the sum relation
$m_M=3m_Lm_R-m_D$ is its other root. The product of the two roots is
$m_Mm_D=m_L^2+m_R^2$, which is the second relation in
\eqref{eq:vieta-relations}, and $(m_L,m_R,m_M)$ satisfies the Markov equation
as well. The triples of the two new pairs, $(m_L,m_M,m_R)$ and
$(m_M,m_R,m_L)$, are permutations of it, and induction on the depth in the
tree proves that every triple $(m_L,m_R,m_D)$ is a Markov triple.

\emph{Coprimality.} If a prime divides two entries of a Markov triple, the
Markov equation shows that it divides the square of the third, hence the
third. The triple of a new pair has the entries $m_L$, $m_R$ and $m_M$, so a
prime dividing two of them would divide $m_L$ and $m_R$, two entries of the
triple of the parent pair. Since the starting triple $(1,2,1)$ has pairwise
coprime entries, induction on the depth proves the last assertion.
\end{proof}

The product relation is the one to lift. Put $\alpha_{0/1}=1$,
$\alpha_{1/1}=1+i$ and $\alpha_{1/0}=1$, and, for each pair $(L,R)$ of the
tree, with mediant $M$ and difference $D$,
\begin{equation}\label{eq:vieta-jump}
 \overline{\alpha_M}=\frac{m_L+i\,m_R}{\alpha_D}.
\end{equation}
The definition proceeds down the tree: when a pair is reached, $\alpha_D$ is
already known, since $D$ is a starting fraction or a mediant formed earlier,
and it remains to see that the division is exact. Three divisibilities
attached to each triangle ensure it. They are exchanged by the circular
permutation $L\to R\to D\to L$: the first says that the division
\eqref{eq:vieta-jump} is exact, and the other two prepare the divisions of the
next generation. The first lemma below carries them from a pair to its two
new pairs; the second runs the induction down the tree.

\begin{lemma}[One step down the tree]\label{lem:vieta-step}
Let $(L,R)$ be a pair of the tree, with mediant $M$ and difference $D$, and
let $\alpha_L$, $\alpha_R$, $\alpha_D$ be Gaussian integers of norms $m_L$, $m_R$,
$m_D$ such that
\begin{equation}\label{eq:three-divisibilities}
 \alpha_D\mid m_L+i\,m_R,\qquad \alpha_L\mid m_R+i\,m_D,\qquad
 \alpha_R\mid m_D+i\,m_L .
\end{equation}
Then the quotient \eqref{eq:vieta-jump} is a Gaussian integer, $N(\alpha_M)=m_M$,
and
\begin{equation}\label{eq:new-divisibilities}
 \alpha_M\mid m_R+i\,m_L,\qquad \alpha_L\mid m_M+i\,m_R,\qquad
 \alpha_R\mid m_L+i\,m_M .
\end{equation}
\end{lemma}

\begin{proof}
\emph{The norm of $\alpha_M$.} The first divisibility in
\eqref{eq:three-divisibilities} makes the quotient \eqref{eq:vieta-jump} a
Gaussian integer, and by \cref{lem:farey-triangle},
$$
 N(\alpha_M)=\frac{N(m_L+i\,m_R)}{N(\alpha_D)}=\frac{m_L^2+m_R^2}{m_D}
 =\frac{m_Mm_D}{m_D}=m_M .
$$

\emph{Divisibility by $\alpha_M$.} Conjugating \eqref{eq:vieta-jump} gives
$\alpha_M\overline{\alpha_D}=m_L-i\,m_R$, and multiplying by $i$,
$$
 \alpha_M\cdot i\,\overline{\alpha_D}=m_R+i\,m_L .
$$

\emph{Divisibility by $\alpha_L$.} We compute modulo $\alpha_L$, which divides
$m_L=\alpha_L\overline{\alpha_L}$. The definition gives
$\overline{\alpha_M}\alpha_D=m_L+i\,m_R\equiv i\,m_R$, and its conjugate gives
$\alpha_M\overline{\alpha_D}=m_L-i\,m_R$. With $m_M=\alpha_M\overline{\alpha_M}$ and
$m_D=\alpha_D\overline{\alpha_D}$ it follows that, modulo $\alpha_L$,
\begin{align*}
 m_M+i\,m_R&\equiv\alpha_M\overline{\alpha_M}+\overline{\alpha_M}\alpha_D
 =\overline{\alpha_M}\,(\alpha_M+\alpha_D),\\
 \overline{\alpha_D}\,(\alpha_M+\alpha_D)&=m_L-i\,m_R+m_D\equiv m_D-i\,m_R
 =-i\,(m_R+i\,m_D).
\end{align*}
The last member is divisible by $\alpha_L$, by the second divisibility in
\eqref{eq:three-divisibilities}. The Gaussian integers $\alpha_L$ and
$\overline{\alpha_D}$ have no common prime factor, because their norms $m_L$ and
$m_D$ are coprime by \cref{lem:farey-triangle}. Unique factorization in $\Z[i]$
therefore gives $\alpha_L\mid\alpha_M+\alpha_D$, and then $\alpha_L\mid m_M+i\,m_R$
by the first line.

\emph{Divisibility by $\alpha_R$.} We compute modulo $\alpha_R$, which divides
$m_R$. The definition gives $\overline{\alpha_M}\alpha_D=m_L+i\,m_R\equiv m_L$,
and therefore, modulo $\alpha_R$,
\begin{align*}
 m_L+i\,m_M&\equiv\overline{\alpha_M}\alpha_D+i\,\alpha_M\overline{\alpha_M}
 =\overline{\alpha_M}\,(\alpha_D+i\,\alpha_M),\\
 \overline{\alpha_D}\,(\alpha_D+i\,\alpha_M)&=m_D+i\,(m_L-i\,m_R)
 =m_D+m_R+i\,m_L\equiv m_D+i\,m_L .
\end{align*}
The last member is divisible by $\alpha_R$, by the third divisibility in
\eqref{eq:three-divisibilities}. The norms $m_R$ and $m_D$ are coprime, so
$\alpha_R$ and $\overline{\alpha_D}$ have no common prime factor, and unique
factorization gives $\alpha_R\mid\alpha_D+i\,\alpha_M$, and then
$\alpha_R\mid m_L+i\,m_M$ by the first line.
\end{proof}

\begin{lemma}[Three divisibilities]\label{lem:three-divisibilities}
Every quotient \eqref{eq:vieta-jump} is a Gaussian integer, and every pair
$(L,R)$ of the tree, with difference $D$, satisfies
\eqref{eq:three-divisibilities} and
\begin{equation}\label{eq:three-norms}
 N(\alpha_L)=m_L,\qquad N(\alpha_R)=m_R,\qquad N(\alpha_D)=m_D.
\end{equation}
\end{lemma}

\begin{proof}
We argue by induction on the depth of the pair in the tree.

\emph{The starting pair.} For $(L,R,D)=(0/1,1/1,1/0)$, the integers $1$, $1+i$
and $1$ have the norms $1$, $2$ and $1$. The first two divisibilities have
the divisor $1$, and the third says that $1+i$ divides
$m_{1/0}+i\,m_{0/1}=1+i$.

\emph{The step.} Assume \eqref{eq:three-norms} and
\eqref{eq:three-divisibilities} for a pair $(L,R)$. By \cref{lem:vieta-step},
the quotient \eqref{eq:vieta-jump} is a Gaussian integer of norm $m_M$, and
\eqref{eq:new-divisibilities} holds. The pair $(L,M)$ has difference $R$, and
its three divisibilities read $\alpha_R\mid m_L+i\,m_M$, $\alpha_L\mid m_M+i\,m_R$
and $\alpha_M\mid m_R+i\,m_L$; the pair $(M,R)$ has difference $L$, and its
three divisibilities read $\alpha_L\mid m_M+i\,m_R$, $\alpha_M\mid m_R+i\,m_L$ and
$\alpha_R\mid m_L+i\,m_M$. Both are the three divisibilities
\eqref{eq:new-divisibilities}, and the norms of the two new pairs are among
$m_L$, $m_R$ and $m_M=N(\alpha_M)$. Both new pairs therefore satisfy
\eqref{eq:three-norms} and \eqref{eq:three-divisibilities}, which completes the
induction.
\end{proof}

The lemma defines $\alpha_M$ for every reduced fraction $M$ of $(0,1)$, since
$M$ is the mediant of exactly one pair of the tree. It remains to identify
this Gaussian integer.

\begin{proposition}[The Vieta jump in the Gaussian integers]\label{prop:vieta-jump}
Let $M=P/Q$ be a reduced fraction of $(0,1)$ with Farey triangle $(L,R,D)$.
Then
\begin{equation}\label{eq:vieta-product}
 m_L+i\,m_R=\overline{\alpha_M}\,\alpha_D,\qquad N(\alpha_M)=m_M,
\end{equation}
and $\alpha_M$ is the occurrence pair of $P/Q$ up to a unit and a conjugation:
if $(s,r)^{\mathsf T}=He$, then $\alpha_M=\varepsilon(s+ir)$ or
$\alpha_M=\varepsilon(s-ir)$ for a unit $\varepsilon$ of $\Z[i]$. In particular, when
$D$ lies in $(0,1)$, the cofactor $\alpha_D$ of $\overline{\alpha_M}$ in
$m_L+i\,m_R$ is the occurrence pair of $D$, up to a unit and a conjugation.
\end{proposition}

\begin{proof}
\Cref{lem:three-divisibilities} shows that \eqref{eq:vieta-jump} is a Gaussian
integer of norm $m_M$, and multiplying \eqref{eq:vieta-jump} by $\alpha_D$
gives the product in \eqref{eq:vieta-product}. The Gaussian integer
$\overline{\alpha_M}$ divides $m_M=\overline{\alpha_M}\alpha_M$ and
$m_L+i\,m_R=\overline{\alpha_M}\alpha_D$, hence every element of the ideal of
$\Z[i]$ generated by $m_M$ and $m_L+i\,m_R$. By \cref{prop:occurrence-in-cube},
this ideal is generated by $s+ir$ or by $s-ir$, so that $\overline{\alpha_M}$
divides one of them. The two sides of this divisibility have the same norm
$m_M$, so their quotient is a Gaussian integer of norm one, a unit:
$\overline{\alpha_M}=\varepsilon'(s\pm ir)$, and conjugation gives
$\alpha_M=\overline{\varepsilon'}(s\mp ir)$. The last assertion is the first
one applied to the fraction $D$, which has its own Farey triangle.
\end{proof}

Up to units and conjugation, the product form is classical in substance: it
follows from the identity $mm'=m_L^2+m_R^2$ and from the characteristic
numbers of Frobenius, both recalled after \cref{prop:occurrence-in-cube} and
applied to $m$ and to $m'=m_D$, together with unique factorization in $\Z[i]$
\cite[\S\S1--2]{Frobenius1913}. For $m_D$ one more step is needed, since
$(m_D,m_L,m_R)$ is not the triple of the Farey parents of $D$: the Farey
triangles with the vertex $D$ are joined by flips that keep $D$, each of which
replaces a triple $(m_D,x,y)$ by $(m_D,y,3m_Dy-x)$, and this leaves the residue
of $x/y$ modulo $m_D$ unchanged up to sign, because $3m_Dy-x\equiv-x$ and
$(x/y)^2\equiv-1\pmod{m_D}$ by the Markov equation. What
\cref{lem:vieta-step,lem:three-divisibilities} add is the exact orientation,
without units, along the whole tree, and it is this orientation that makes the
recursion possible.

The proposition turns the computation of occurrence pairs into a descent.
To obtain the occurrence pair of $P/Q$, one follows the path of $P/Q$ in the
Farey tree and keeps only the three Gaussian integers of the current
triangle; each step costs one exact division in $\Z[i]$.
\Cref{fig:vieta-tree} shows the first four generations. For the triangle
$L=1/2$, $R=2/3$, $D=1/1$, whose labels are $5$, $29$ and $2$, the rule
\eqref{eq:vieta-jump} gives
$$
 \overline{\alpha_{3/5}}=\frac{5+29\,i}{1+i}=\frac{(5+29\,i)(1-i)}2
 =\frac{34+24\,i}2=17+12\,i ,
$$
and indeed $17^2+12^2=433=m_{3/5}$. Along the fractions $P/(P+1)$, whose right
parent is $1/1$ and whose labels are Pell numbers, the rule gives
successively
\begin{align*}
 29+2i&=(5+12i)(1-2i),\\
 169+2i&=(29+12i)(5-2i),\\
 985+2i&=(29+70i)(5-12i),
\end{align*}
that is $169=5^2+12^2$, $985=29^2+12^2$ and $5741=29^2+70^2$, each second
factor being the Gaussian integer found two generations earlier. Along the
fractions $1/Q$ it recovers the identity $1+233i=(21+13i)(5+8i)$ of the
label $610$, with both factors multiplied by $-1$.

\begin{figure}[tb]
\centering
% Authors: Dr. Denys Dutykh (Mathematics Department, Khalifa University of Science and
% Technology, Abu Dhabi, UAE) and Prof. Laurent Vuillon (Univ. Savoie Mont Blanc, CNRS, LAMA)
% Generated by supplement/code/make_figures.py from exact data; do not edit by hand.
\begin{tikzpicture}[font=\scriptsize]
\node[draw=slate, line width=0.5pt, dotted, fill=white, rounded corners=1.5pt, inner sep=2pt, align=center] (f0x1) at (0,0) {$0/1$\\$1$\\$1$};
\node[draw=wine, line width=0.9pt, fill=white, rounded corners=1.5pt, inner sep=2pt, align=center] (f1x1) at (13.44,0) {$D=1/1$\\$2$\\$1+i$};
\node[draw=slate, line width=0.5pt, dotted, fill=white, rounded corners=1.5pt, inner sep=2pt, align=center] (f1x0) at (6.72,0) {$1/0$\\$1$\\$1$};
\node[draw=wine, line width=0.9pt, fill=white, rounded corners=1.5pt, inner sep=2pt, align=center] (f1x2) at (6.72,-1.42) {$L=1/2$\\$5$\\$1-2i$};
\node[draw=mist, line width=0.5pt, fill=white, rounded corners=1.5pt, inner sep=2pt, align=center] (f1x3) at (3.36,-2.84) {$1/3$\\$13$\\$3-2i$};
\node[draw=wine, line width=0.9pt, fill=white, rounded corners=1.5pt, inner sep=2pt, align=center] (f2x3) at (10.08,-2.84) {$R=2/3$\\$29$\\$5-2i$};
\node[draw=mist, line width=0.5pt, fill=white, rounded corners=1.5pt, inner sep=2pt, align=center] (f1x4) at (1.68,-4.26) {$1/4$\\$34$\\$-5-3i$};
\node[draw=mist, line width=0.5pt, fill=white, rounded corners=1.5pt, inner sep=2pt, align=center] (f2x5) at (5.04,-4.26) {$2/5$\\$194$\\$13-5i$};
\node[draw=navy, line width=1.1pt, fill=sand, rounded corners=1.5pt, inner sep=2pt, align=center] (f3x5) at (8.4,-4.26) {$M=3/5$\\$433$\\$17-12i$};
\node[draw=mist, line width=0.5pt, fill=white, rounded corners=1.5pt, inner sep=2pt, align=center] (f3x4) at (11.76,-4.26) {$3/4$\\$169$\\$5-12i$};
\node[draw=mist, line width=0.5pt, fill=white, rounded corners=1.5pt, inner sep=2pt, align=center] (f1x5) at (0.84,-5.68) {$1/5$\\$89$\\$-5-8i$};
\node[draw=mist, line width=0.5pt, fill=white, rounded corners=1.5pt, inner sep=2pt, align=center] (f2x7) at (2.52,-5.68) {$2/7$\\$1325$\\$34-13i$};
\node[draw=mist, line width=0.5pt, fill=white, rounded corners=1.5pt, inner sep=2pt, align=center] (f3x8) at (4.2,-5.68) {$3/8$\\$7561$\\$-75-44i$};
\node[draw=mist, line width=0.5pt, fill=white, rounded corners=1.5pt, inner sep=2pt, align=center] (f3x7) at (5.88,-5.68) {$3/7$\\$2897$\\$44-31i$};
\node[draw=mist, line width=0.5pt, fill=white, rounded corners=1.5pt, inner sep=2pt, align=center] (f4x7) at (7.56,-5.68) {$4/7$\\$6466$\\$-29-75i$};
\node[draw=mist, line width=0.5pt, fill=white, rounded corners=1.5pt, inner sep=2pt, align=center] (f5x8) at (9.24,-5.68) {$5/8$\\$37666$\\$75-179i$};
\node[draw=mist, line width=0.5pt, fill=white, rounded corners=1.5pt, inner sep=2pt, align=center] (f5x7) at (10.92,-5.68) {$5/7$\\$14701$\\$99-70i$};
\node[draw=mist, line width=0.5pt, fill=white, rounded corners=1.5pt, inner sep=2pt, align=center] (f4x5) at (12.6,-5.68) {$4/5$\\$985$\\$29-12i$};
\begin{scope}[on background layer]
\draw[mist, line width=0.6pt] (f0x1) -- (f1x2);
\draw[mist, line width=0.6pt] (f1x2) -- (f1x3);
\draw[mist, line width=0.6pt] (f1x3) -- (f1x4);
\draw[mist, line width=0.6pt] (f1x3) -- (f2x5);
\draw[mist, line width=0.6pt] (f1x4) -- (f1x5);
\draw[mist, line width=0.6pt] (f1x4) -- (f2x7);
\draw[mist, line width=0.6pt] (f2x3) -- (f3x4);
\draw[mist, line width=0.6pt] (f2x5) -- (f3x7);
\draw[mist, line width=0.6pt] (f2x5) -- (f3x8);
\draw[mist, line width=0.6pt] (f3x4) -- (f4x5);
\draw[mist, line width=0.6pt] (f3x4) -- (f5x7);
\draw[mist, line width=0.6pt] (f3x5) -- (f4x7);
\draw[mist, line width=0.6pt] (f3x5) -- (f5x8);
\draw[wine, line width=0.8pt, dash pattern=on 3pt off 2pt] (f1x2) -- (f2x3);
\draw[wine, line width=0.8pt, dash pattern=on 3pt off 2pt] (f2x3) -- (f1x1);
\draw[wine, line width=0.8pt, dash pattern=on 3pt off 2pt] (f1x1) -- (f1x2);
\draw[navy, line width=1.0pt] (f1x2) -- (f3x5);
\draw[navy, line width=1.0pt] (f2x3) -- (f3x5);
\end{scope}
\end{tikzpicture}

\caption{The first four generations of the Farey tree, with the label $m$
and the Gaussian integer $\alpha$ of each fraction, computed by the rule
\eqref{eq:vieta-jump} alone from the values at the three starting fractions
$0/1$, $1/1$ and $1/0$, drawn at the top. Each fraction is joined to the younger of its
two Farey parents, and $1/2$ to both. In wine and dashed, the Farey triangle
$L=1/2$, $R=2/3$, $D=1/1$; in navy, the edges from its parents to the
mediant $M=3/5$, whose Gaussian integer $17-12i$ is the conjugate of the
quotient of $5+29i$ by~$1+i$.}
\label{fig:vieta-tree}
\end{figure}

\begin{remark}[The pair determines the occurrence]\label{rem:pair-determines}
Two distinct occurrences of one label have distinct occurrence pairs, even
up to units and conjugation: equal pairs have equal costs, and
\cref{lem:fixed-cost-injectivity} applies. \Cref{prop:vieta-jump} computes the
occurrence pair from the Farey triangle, that is, from the position of the
occurrence in the Markov tree, and not from its label. A rule that produced
the occurrence pair from the label alone, correctly for every occurrence,
would therefore prove the uniqueness conjecture; \cref{sec:directives}
states a conjecture of this kind.
\end{remark}

The finite part of the proof, to which we now turn, needs no recursion: it
computes the whole fiber of each label, in which the occurrence pair is the
one given by the Gram factor of \cref{prop:occurrence-cost}.

\subsection{An elementary primality criterion}

The factorizations used below rest on the following classical criterion,
which we state for completeness. Suppose $n>2$, the complete factorization of $n-1$ into proved
primes is supplied, and an integer $a$ satisfies
\begin{equation}\label{eq:full-lucas}
 a^{n-1}\equiv1\pmod n,\qquad
 \gcd\bigl(a^{(n-1)/r}-1,n\bigr)=1
 \quad\text{for every prime }r\mid n-1.
\end{equation}
Then $n$ is prime. Indeed, for any prime divisor $s$ of $n$, these
conditions make the order of $a$ modulo $s$ exactly $n-1$. Since that
order divides $s-1$, one has $s\geqslant n$, hence $s=n$.

Applied recursively, with the primes of $n-1$ themselves established the
same way, the criterion terminates at the prime two and establishes
primality by integer arithmetic alone: no probabilistic test and no search
for factors enters the verification, only the confirmation of factors
already proposed.

Each label is computed twice and the two values are compared: once from
the floor increments of the lower Christoffel word, and once from the Cohn
matrix built by Stern--Brocot mediants from the two boundary matrices. The
occurrence pair is computed from the central Gram factor $H$ of
\cref{prop:occurrence-cost}. For each label one then enumerates the
$2^t$ sign choices in \eqref{eq:gaussian-sign-cube} and passes to the
quotient by conjugation; that every normalized pair is obtained exactly
twice, as \cref{prop:gaussian-fibre} predicts, is itself a check on the
enumeration. The Euclidean algorithm then supplies the quotient sequence
and the cost of every pair.

\subsection{The finite windows}

The finite part of the proof concerns the reduced fractions listed in
\cref{tab:gaussian-windows}: in row $P$, for $8\leqslant P\leqslant13$, every
denominator from $P+1$ up to the last denominator before the junction
$Q_P$ of \cref{sec:assembly}. For each of these $248$ labels, the prime
factorization, the primality of its factors by the recursive use of
\eqref{eq:full-lucas}, the splitting of each prime and the costs of the
resulting pairs are finitely many assertions about integers. They are
recorded as part~\ref{it:fv-windows} of the finite verifications,
\cref{prop:finite-verifications}, and \cref{prop:finite-windows} draws
their consequence: in every finite window the occurrence pair is the unique
representation of minimum cost.

\begin{table}[tb]
\centering
\caption{The complete finite part of the minimality proof. Each row includes
all $Q$ in the displayed interval with $\gcd(P,Q)=1$. The pair count includes
every normalized primitive positive representation of every source, and the
last column gives the largest individual fiber. The later sections cover the
unbounded denominators beyond these windows.}
\label{tab:gaussian-windows}
\small
\setlength{\tabcolsep}{11pt}
\renewcommand{\arraystretch}{1.12}
\begin{tabular}{@{}crrrr@{}}
\toprule
$P$ & Denominators $Q$ & Sources & Pairs & Largest fiber \\
\midrule
8 & $9$--$28$ & 10 & 22 & 4 \\
9 & $10$--$34$ & 17 & 56 & 8 \\
10 & $11$--$38$ & 11 & 33 & 8 \\
11 & $12$--$75$ & 59 & 324 & 32 \\
12 & $13$--$100$ & 29 & 211 & 32 \\
13 & $14$--$145$ & 122 & 992 & 64 \\
\midrule
Total & & 248 & 1638 & 64 \\
\bottomrule
\end{tabular}
\end{table}

The statement on the finite windows is stronger than weak minimality, since
it also asserts uniqueness of the minimizing representation. The unbounded
part requires only weak minimality. It uses the finite windows as exact
boundary data and separately excludes every lower-cost competitor on each
infinite tail.
